% Enunciado: Seções 1.1 e 1.2 (o problema é dado; apresentar de forma breve)
% Tamanho sugerido: ~2/3 de coluna

\section{Introdução}
\label{sec:introducao}

% Parágrafo 1 -- contexto: concessionária de energia, uma única equipe que
% parte da base e atende um ativo por vez; cada ativo j tem posição
% (x_j, y_j), duração p_j, prazo d_j e criticidade w_j.

% Formulação compacta, definindo a notação usada no resto do texto:
% - tempo de deslocamento t_ij arredondado (\label{eq:deslocamento});
% - instante de conclusão C_{pi_k} (\label{eq:conclusao});
% - função de avaliação h(pi) = sum w_j max(0, C_j - d_j) (\label{eq:fitness}).
% Uma frase sobre o tamanho do espaço (n!, n = 150 => >260 algarismos), o
% fato de toda permutação ser viável e a tensão entre ir ao mais próximo e
% ir ao mais urgente.

% Parágrafo 2 -- objetivo: comparar busca local (greedy descent com random
% restart, simulated annealing) e uma metaheurística populacional
% (algoritmo genético) sob o mesmo orçamento de avaliações, observando o
% equilíbrio exploração x intensificação. Citar \cite{poole2010}.

% Parágrafo 3 -- organização do texto, referenciando:
% Seção~\ref{sec:metodologia}, Seção~\ref{sec:algoritmos},
% Seção~\ref{sec:resultados} (parte central) e Seção~\ref{sec:conclusao}.
